\section{Proofs of Estimators (Multivariate GMM - Classification)}

We now consider the multivariate case, where each observation has
$p$ features.

For $p>1$, we consider a \hl{multivariate Gaussian mixture model}:
\begin{equation*}
X\mid Z=k
\sim
\mathcal{N}_p(\boldsymbol{\mu}_k,\boldsymbol{\Sigma}_k),
\qquad
k=1,\ldots,K,
\end{equation*}
where
\begin{equation*}
\boldsymbol{\mu}_k\in\mathbb{R}^p
\end{equation*}
is the mean vector of group $k$ and
\begin{equation*}
\boldsymbol{\Sigma}_k\in\mathbb{R}^{p\times p}
\end{equation*}
is the covariance matrix of group $k$.

Since $\boldsymbol{\Sigma}_k$ is a covariance matrix, it is symmetric
positive definite:
\begin{equation*}
\boldsymbol{\Sigma}_k^\top
=
\boldsymbol{\Sigma}_k,
\qquad
\boldsymbol{\Sigma}_k\succ0.
\end{equation*}

The one-hot representation of the group variable is
\begin{equation*}
\tilde{Z}\sim\mathcal{M}(1,p_1,\ldots,p_K),
\end{equation*}
where
\begin{equation*}
p_k=P(Z=k),
\qquad
\sum_{k=1}^Kp_k=1.
\end{equation*}

The parameter vector is now
\begin{equation*}
\theta=
\left(
p_1,\ldots,p_K,
\boldsymbol{\mu}_1,\ldots,\boldsymbol{\mu}_K,
\boldsymbol{\Sigma}_1,\ldots,\boldsymbol{\Sigma}_K
\right).
\end{equation*}

For group $k$, the multivariate Gaussian density is
\begin{equation*}
f_k(\boldsymbol{x})
=
\frac{1}
{(2\pi)^{p/2}
|\boldsymbol{\Sigma}_k|^{1/2}}
\exp\left(
-\frac{1}{2}
(\boldsymbol{x}-\boldsymbol{\mu}_k)^\top
\boldsymbol{\Sigma}_k^{-1}
(\boldsymbol{x}-\boldsymbol{\mu}_k)
\right).
\end{equation*}

We observe
\begin{equation*}
\underline{X}
=
(\boldsymbol{X}_1,\ldots,\boldsymbol{X}_n),
\end{equation*}
where
\begin{equation*}
\boldsymbol{X}_i\in\mathbb{R}^p,
\end{equation*}
and denote the corresponding group variables by
\begin{equation*}
\underline{Z}=(Z_1,\ldots,Z_n).
\end{equation*}

Using the one-hot representation $\tilde{Z}_{ik}$, where
\begin{equation*}
\tilde{Z}_{ik}
=
\begin{cases}
1 & \text{if } Z_i=k,\\
0 & \text{otherwise},
\end{cases}
\end{equation*}
the complete-data log-likelihood is
\begin{equation*}
\ln L(\underline{X},Z,\theta)
=
\sum_{i=1}^{n}\sum_{k=1}^{K}
\tilde{Z}_{ik}
\left[
\ln p_k
-\frac{p}{2}\ln(2\pi)
-\frac{1}{2}\ln|\boldsymbol{\Sigma}_k|
-\frac{1}{2}
(\boldsymbol{x}_i-\boldsymbol{\mu}_k)^\top
\boldsymbol{\Sigma}_k^{-1}
(\boldsymbol{x}_i-\boldsymbol{\mu}_k)
\right].
\end{equation*}

We now derive the maximum likelihood estimators of the three types
of parameters:
\begin{equation*}
p_k,\qquad
\boldsymbol{\mu}_k,\qquad
\boldsymbol{\Sigma}_k.
\end{equation*}

As in the univariate case, we treat the group memberships
$Z_1,\ldots,Z_n$ as observed. Therefore, these are the
\hl{complete-data maximum likelihood estimators}.


%------------------------------------------------------------
\subsection{Estimation of the Mixing Proportions $p_k$}
%------------------------------------------------------------

The mixing proportions satisfy the constraint
\begin{equation*}
\sum_{k=1}^Kp_k=1.
\end{equation*}

We therefore introduce a Lagrange multiplier $\lambda$ and define
the Lagrangian
\begin{equation*}
F(\theta,\lambda)
=
\ln L(\underline{X},Z,\theta)
+
\lambda
\left(
1-\sum_{k=1}^Kp_k
\right).
\end{equation*}

We use the conditions
\begin{equation*}
\frac{\partial F}{\partial p_k}=0,
\qquad
\frac{\partial F}{\partial\lambda}=0.
\end{equation*}

Define
\begin{equation*}
n_k
=
\sum_{i=1}^n\tilde{Z}_{ik},
\end{equation*}
where $n_k$ is the number of observations belonging to group $k$.

Keeping only the terms that depend on the mixing proportions,
the Lagrangian becomes
\begin{equation*}
F
=
\sum_{k=1}^K n_k\ln p_k
+
\lambda
\left(
1-\sum_{k=1}^Kp_k
\right)
+
\text{constant}.
\end{equation*}

\paragraph{Step 1: Differentiate with respect to $p_k$.}

For a fixed group $k$,
\begin{equation*}
\frac{\partial F}{\partial p_k}
=
\frac{n_k}{p_k}-\lambda.
\end{equation*}

At the optimum,
\begin{equation*}
\frac{n_k}{p_k}-\lambda=0.
\end{equation*}

Therefore,
\begin{equation*}
p_k=\frac{n_k}{\lambda}.
\end{equation*}

\paragraph{Step 2: Differentiate with respect to $\lambda$.}

The derivative with respect to the Lagrange multiplier gives back
the constraint:
\begin{equation*}
\frac{\partial F}{\partial\lambda}
=
1-\sum_{k=1}^Kp_k
=
0.
\end{equation*}

Therefore,
\begin{equation*}
\sum_{k=1}^Kp_k=1.
\end{equation*}

Substituting
\begin{equation*}
p_k=\frac{n_k}{\lambda}
\end{equation*}
into the constraint gives
\begin{equation*}
\sum_{k=1}^K\frac{n_k}{\lambda}=1.
\end{equation*}

Since every observation belongs to exactly one group,
\begin{equation*}
\sum_{k=1}^Kn_k=n.
\end{equation*}

Therefore,
\begin{equation*}
\frac{n}{\lambda}=1
\qquad\Longrightarrow\qquad
\lambda=n.
\end{equation*}

Hence,
\begin{equation*}
\boxed{
\hat{p}_k=\frac{n_k}{n}
}.
\end{equation*}

Since
\begin{equation*}
n_k=\sum_{i=1}^n\tilde{Z}_{ik},
\end{equation*}
we can equivalently write
\begin{equation*}
\boxed{
\hat{p}_k
=
\frac{1}{n}
\sum_{i=1}^n\tilde{Z}_{ik}
}.
\end{equation*}

Thus, exactly as in the univariate case, the estimated mixing
proportion of group $k$ is the \hl{proportion of observations
belonging to group $k$}.


%------------------------------------------------------------
\subsection{Estimation of the Mean Vector}
%------------------------------------------------------------

We now estimate the mean vector
\begin{equation*}
\boldsymbol{\mu}_k\in\mathbb{R}^p
\end{equation*}
of group $k$.

Unlike the mixing proportions, the mean vectors are not subject to
a constraint. Therefore, we directly differentiate the
log-likelihood with respect to $\boldsymbol{\mu}_k$.

Keeping only the terms that depend on $\boldsymbol{\mu}_k$ gives
\begin{equation*}
\ln L
=
-\frac{1}{2}
\sum_{i=1}^n
\tilde{Z}_{ik}
(\boldsymbol{x}_i-\boldsymbol{\mu}_k)^\top
\boldsymbol{\Sigma}_k^{-1}
(\boldsymbol{x}_i-\boldsymbol{\mu}_k)
+
\text{constant}.
\end{equation*}

Since $\boldsymbol{\Sigma}_k$ is symmetric positive definite,
its inverse is also symmetric:
\begin{equation*}
(\boldsymbol{\Sigma}_k^{-1})^\top
=
\boldsymbol{\Sigma}_k^{-1}.
\end{equation*}

We use the standard matrix derivative
\begin{equation*}
\frac{\partial}{\partial\boldsymbol{\mu}}
\left[
(\boldsymbol{x}-\boldsymbol{\mu})^\top
A
(\boldsymbol{x}-\boldsymbol{\mu})
\right]
=
-2A(\boldsymbol{x}-\boldsymbol{\mu}),
\end{equation*}
when $A$ is symmetric.

Therefore,
\begin{align*}
\frac{\partial\ln L}
{\partial\boldsymbol{\mu}_k}
&=
-\frac{1}{2}
\sum_{i=1}^n
\tilde{Z}_{ik}
\left[
-2\boldsymbol{\Sigma}_k^{-1}
(\boldsymbol{x}_i-\boldsymbol{\mu}_k)
\right]\\
&=
\boldsymbol{\Sigma}_k^{-1}
\sum_{i=1}^n
\tilde{Z}_{ik}
(\boldsymbol{x}_i-\boldsymbol{\mu}_k).
\end{align*}

At the maximum likelihood estimator,
\begin{equation*}
\frac{\partial\ln L}
{\partial\boldsymbol{\mu}_k}
=
\boldsymbol{0}.
\end{equation*}

Thus,
\begin{equation*}
\boldsymbol{\Sigma}_k^{-1}
\sum_{i=1}^n
\tilde{Z}_{ik}
(\boldsymbol{x}_i-\boldsymbol{\mu}_k)
=
\boldsymbol{0}.
\end{equation*}

Since $\boldsymbol{\Sigma}_k$ is invertible, we can multiply both
sides by $\boldsymbol{\Sigma}_k$:
\begin{equation*}
\sum_{i=1}^n
\tilde{Z}_{ik}
(\boldsymbol{x}_i-\boldsymbol{\mu}_k)
=
\boldsymbol{0}.
\end{equation*}

Expanding the expression gives
\begin{equation*}
\sum_{i=1}^n
\tilde{Z}_{ik}\boldsymbol{x}_i
-
\boldsymbol{\mu}_k
\sum_{i=1}^n\tilde{Z}_{ik}
=
\boldsymbol{0}.
\end{equation*}

Using
\begin{equation*}
n_k=\sum_{i=1}^n\tilde{Z}_{ik},
\end{equation*}
we obtain
\begin{equation*}
\sum_{i=1}^n
\tilde{Z}_{ik}\boldsymbol{x}_i
-
n_k\boldsymbol{\mu}_k
=
\boldsymbol{0}.
\end{equation*}

Therefore,
\begin{equation*}
n_k\boldsymbol{\mu}_k
=
\sum_{i=1}^n
\tilde{Z}_{ik}\boldsymbol{x}_i
\end{equation*}
and hence
\begin{equation*}
\boxed{
\hat{\boldsymbol{\mu}}_k
=
\frac{1}{n_k}
\sum_{i=1}^n
\tilde{Z}_{ik}\boldsymbol{x}_i
}.
\end{equation*}

Thus, the maximum likelihood estimator of the mean vector is the
\hl{sample mean vector of the observations belonging to group $k$}.


%------------------------------------------------------------
\subsection{Estimation of the Covariance Matrix}
%------------------------------------------------------------

Finally, we estimate the covariance matrix
\begin{equation*}
\boldsymbol{\Sigma}_k.
\end{equation*}

For a fixed group $k$, the terms of the log-likelihood that depend
on $\boldsymbol{\Sigma}_k$ are
\begin{equation*}
\ln L
=
\sum_{i=1}^n
\tilde{Z}_{ik}
\left[
-\frac{1}{2}\ln|\boldsymbol{\Sigma}_k|
-\frac{1}{2}
(\boldsymbol{x}_i-\boldsymbol{\mu}_k)^\top
\boldsymbol{\Sigma}_k^{-1}
(\boldsymbol{x}_i-\boldsymbol{\mu}_k)
\right]
+
\text{constant}.
\end{equation*}

For convenience, define
\begin{equation*}
\boldsymbol{d}_{ik}
=
\boldsymbol{x}_i-\boldsymbol{\mu}_k.
\end{equation*}

Then the relevant part of the log-likelihood can be written as
\begin{equation*}
\ln L
=
-\frac{1}{2}
\sum_{i=1}^n
\tilde{Z}_{ik}
\left[
\ln|\boldsymbol{\Sigma}_k|
+
\boldsymbol{d}_{ik}^\top
\boldsymbol{\Sigma}_k^{-1}
\boldsymbol{d}_{ik}
\right]
+
\text{constant}.
\end{equation*}

We use the following standard matrix derivatives:
\begin{equation*}
\frac{\partial}{\partial\boldsymbol{\Sigma}}
\ln|\boldsymbol{\Sigma}|
=
\boldsymbol{\Sigma}^{-1},
\end{equation*}
and
\begin{equation*}
\frac{\partial}{\partial\boldsymbol{\Sigma}}
\left(
\boldsymbol{d}^\top
\boldsymbol{\Sigma}^{-1}
\boldsymbol{d}
\right)
=
-\boldsymbol{\Sigma}^{-1}
\boldsymbol{d}\boldsymbol{d}^\top
\boldsymbol{\Sigma}^{-1},
\end{equation*}
where $\boldsymbol{\Sigma}$ is symmetric positive definite.

Therefore,
\begin{align*}
\frac{\partial\ln L}
{\partial\boldsymbol{\Sigma}_k}
&=
-\frac{1}{2}
\sum_{i=1}^n
\tilde{Z}_{ik}
\left[
\boldsymbol{\Sigma}_k^{-1}
-
\boldsymbol{\Sigma}_k^{-1}
\boldsymbol{d}_{ik}\boldsymbol{d}_{ik}^\top
\boldsymbol{\Sigma}_k^{-1}
\right]\\
&=
-\frac{1}{2}
\sum_{i=1}^n
\tilde{Z}_{ik}
\left[
\boldsymbol{\Sigma}_k^{-1}
-
\boldsymbol{\Sigma}_k^{-1}
(\boldsymbol{x}_i-\boldsymbol{\mu}_k)
(\boldsymbol{x}_i-\boldsymbol{\mu}_k)^\top
\boldsymbol{\Sigma}_k^{-1}
\right].
\end{align*}

At the maximum likelihood estimator,
\begin{equation*}
\frac{\partial\ln L}
{\partial\boldsymbol{\Sigma}_k}
=
\boldsymbol{0}.
\end{equation*}

Therefore,
\begin{equation*}
\sum_{i=1}^n
\tilde{Z}_{ik}
\left[
\boldsymbol{\Sigma}_k^{-1}
-
\boldsymbol{\Sigma}_k^{-1}
(\boldsymbol{x}_i-\boldsymbol{\mu}_k)
(\boldsymbol{x}_i-\boldsymbol{\mu}_k)^\top
\boldsymbol{\Sigma}_k^{-1}
\right]
=
\boldsymbol{0}.
\end{equation*}

Expanding the sum gives
\begin{equation*}
\left(
\sum_{i=1}^n\tilde{Z}_{ik}
\right)
\boldsymbol{\Sigma}_k^{-1}
-
\boldsymbol{\Sigma}_k^{-1}
\left[
\sum_{i=1}^n
\tilde{Z}_{ik}
(\boldsymbol{x}_i-\boldsymbol{\mu}_k)
(\boldsymbol{x}_i-\boldsymbol{\mu}_k)^\top
\right]
\boldsymbol{\Sigma}_k^{-1}
=
\boldsymbol{0}.
\end{equation*}

Using
\begin{equation*}
n_k=\sum_{i=1}^n\tilde{Z}_{ik},
\end{equation*}
we obtain
\begin{equation*}
n_k\boldsymbol{\Sigma}_k^{-1}
-
\boldsymbol{\Sigma}_k^{-1}
\left[
\sum_{i=1}^n
\tilde{Z}_{ik}
(\boldsymbol{x}_i-\boldsymbol{\mu}_k)
(\boldsymbol{x}_i-\boldsymbol{\mu}_k)^\top
\right]
\boldsymbol{\Sigma}_k^{-1}
=
\boldsymbol{0}.
\end{equation*}

Multiplying on the left and right by
$\boldsymbol{\Sigma}_k$ gives
\begin{equation*}
n_k\boldsymbol{\Sigma}_k
-
\sum_{i=1}^n
\tilde{Z}_{ik}
(\boldsymbol{x}_i-\boldsymbol{\mu}_k)
(\boldsymbol{x}_i-\boldsymbol{\mu}_k)^\top
=
\boldsymbol{0}.
\end{equation*}

Therefore,
\begin{equation*}
\sum_{i=1}^n
\tilde{Z}_{ik}
(\boldsymbol{x}_i-\boldsymbol{\mu}_k)
(\boldsymbol{x}_i-\boldsymbol{\mu}_k)^\top
=
n_k\boldsymbol{\Sigma}_k.
\end{equation*}

Hence,
\begin{equation*}
\boldsymbol{\Sigma}_k
=
\frac{1}{n_k}
\sum_{i=1}^n
\tilde{Z}_{ik}
(\boldsymbol{x}_i-\boldsymbol{\mu}_k)
(\boldsymbol{x}_i-\boldsymbol{\mu}_k)^\top.
\end{equation*}

Since the maximum likelihood estimator of the mean vector is
$\hat{\boldsymbol{\mu}}_k$, the covariance estimator is therefore
\begin{equation*}
\boxed{
\hat{\boldsymbol{\Sigma}}_k
=
\frac{1}{n_k}
\sum_{i=1}^n
\tilde{Z}_{ik}
(\boldsymbol{x}_i-\hat{\boldsymbol{\mu}}_k)
(\boldsymbol{x}_i-\hat{\boldsymbol{\mu}}_k)^\top
}.
\end{equation*}

Thus, the maximum likelihood estimator of the covariance matrix is
the \hl{empirical covariance matrix of the observations belonging
to group $k$}.

Notice that, as in the univariate case, the denominator is $n_k$,
not $n_k-1$. This is because we are deriving the maximum likelihood
estimator rather than the unbiased estimator of the covariance
matrix.


%------------------------------------------------------------
\subsubsection*{Summary of the Three Estimators}
%------------------------------------------------------------

For a multivariate Gaussian mixture with known group memberships,
the maximum likelihood estimators are
\begin{equation*}
\boxed{
\hat{p}_k
=
\frac{n_k}{n}
}
\end{equation*}

\begin{equation*}
\boxed{
\hat{\boldsymbol{\mu}}_k
=
\frac{1}{n_k}
\sum_{i=1}^n
\tilde{Z}_{ik}\boldsymbol{x}_i
}
\end{equation*}

and
\begin{equation*}
\boxed{
\hat{\boldsymbol{\Sigma}}_k
=
\frac{1}{n_k}
\sum_{i=1}^n
\tilde{Z}_{ik}
(\boldsymbol{x}_i-\hat{\boldsymbol{\mu}}_k)
(\boldsymbol{x}_i-\hat{\boldsymbol{\mu}}_k)^\top
}.
\end{equation*}

In other words,
\begin{equation*}
\boxed{
\begin{aligned}
\hat{p}_k
&=\text{proportion of observations in group }k,\\[2mm]
\hat{\boldsymbol{\mu}}_k
&=\text{mean vector of the observations in group }k,\\[2mm]
\hat{\boldsymbol{\Sigma}}_k
&=\text{covariance matrix of the observations in group }k.
\end{aligned}
}
\end{equation*}

These are exactly the multivariate empirical estimators applied
separately within each group.

The transition from the univariate to the multivariate case is
therefore straightforward:
\begin{equation*}
\boxed{
\begin{aligned}
p_k
&\longrightarrow p_k,\\
\mu_k
&\longrightarrow \boldsymbol{\mu}_k,\\
\sigma_k^2
&\longrightarrow \boldsymbol{\Sigma}_k.
\end{aligned}
}
\end{equation*}

The scalar squared deviation
\begin{equation*}
(x_i-\mu_k)^2
\end{equation*}
is replaced by the outer product
\begin{equation*}
(\boldsymbol{x}_i-\boldsymbol{\mu}_k)
(\boldsymbol{x}_i-\boldsymbol{\mu}_k)^\top,
\end{equation*}
which produces the $p\times p$ covariance matrix.